\documentclass[10pt,a4paper]{amsart}
\usepackage[margin=19mm]{geometry}
\usepackage{amsmath,amssymb,mathtools}
\usepackage[scr=rsfs,cal=euler]{mathalfa}
\usepackage{xfrac}
\usepackage[unicode,hidelinks,bookmarks=false]{hyperref}
\newcommand{\ie}{i.e.\ }
\newcommand{\cf}{cf.\ }
\newcommand{\resp}{respectively\ }
\newcommand{\Subsection}{\subsection}
\providecommand{\nobreakdash}{\nobreak\hbox{-}}
\setlength{\emergencystretch}{2em}
\multlinegap0pt
\usepackage{amssymb}
\usepackage{mathtools}
\usepackage{xcolor}
\usepackage[shortlabels]{enumitem}
\newtheorem{lemma}{Lemma}
\newtheorem{proposition}{Proposition}
\usepackage{cleveref}
\crefname{lemma}{Lemma}{Lemmas}
\crefname{proposition}{Proposition}{Propositions}
\newcommand{\Res}{\operatorname{Res}}
\newcommand{\Tr}{\operatorname{Tr}}
\newcommand{\Z}{\mathbf{Z}}
\newcommand{\dd}{\mathrm{d}}
\title{Split-prime supercongruence at the mixed CM point $(\tfrac16,\tfrac13;1)$}
\author{Alex Shvets}
\date{Source: arXiv:2605.19773v1 (CC BY 4.0)}
\begin{document}
\maketitle

\noindent\emph{Source and licence.} Split-prime supercongruence at the mixed CM point $(\tfrac16,\tfrac13;1)$. Version arXiv:2605.19773v1 (CC BY 4.0), distributed by arXiv under the Creative Commons Attribution 4.0 International licence, http://creativecommons.org/licenses/by/4.0/, as recorded in the arXiv licence metadata for this exact version and shown on the abstract page https://arxiv.org/abs/2605.19773v1. Full licence notice: \texttt{licenses/arxiv-2605.19773-CC-BY-4.0.txt}.\par\smallskip

\noindent\textit{Editorial scope.} Counted unit: the article's proposition prop:pole-estimate (source0.tex lines 745--750). The article's complete original proof of it is delivered with it (source0.tex lines 752--794, one \texttt{proof} environment of 43 lines). No mathematics is written, repaired or re-derived: the statement and the proof are the article's own lines, in the article's own notation, and no retained line is substituted or re-flowed.  Retained uncounted as the source-local context the proof needs: lines 707--713.  Nothing was omitted from any retained span, so the package carries no omission record.  No retained line contains a citation, so the wrapper carries no bibliography and no bibliography entry.  No retained line contains a figure environment, an image-inclusion command or drawing code, so no image is delivered and the package declares no asset dependency.  Editorial formatting only, in addition: an AMS \texttt{amsart} preamble that restates the article's own declarations and macros used by the retained lines, namely the environments \texttt{lemma}, \texttt{proposition} on separate counters, together with the standard packages those lines need.  The retained environments print in this document as Lemma 1, Proposition 1 in order of appearance, whereas the article prints the same items with the numbering of its own full document.  The complete native source of the article as distributed in the arXiv e-print for this exact version is bundled unchanged as \texttt{sources/200940/source0.tex}.
\begin{lemma}[Trace-residue formula on the tame $\mu_3$-chart]\label{lem:trace-residue-tame}
Let $A=\Z_p[\zeta_3][[r]]$, let $F(r)\in r\,A$ be a Frobenius lift on $A$, i.e.\ $F(r)\equiv r^p\pmod p$. Thus $F$ is distinguished of Weierstrass degree $p$; set $R=F(r)$, so $B:=\Z_p[\zeta_3][[R]]\subset A$ is a finite flat extension of rank $p$. Let $J_\Phi(r)\in A^\times$ denote the Jacobian of $\Phi$ with respect to the chosen weight-$5$ trivialization $e$, so that the normalized weight-$5$ Katz trace satisfies $\mathcal C_{\Phi,5}(he)=\bigl(p^{-1}\Tr_{A/B}(h\,J_\Phi^{-1})\bigr)\,e$ for $h$ in the field of fractions of $A$. Then for every $K\ge 1$ and every $h\in\operatorname{Frac}(A)$ with at worst a Laurent principal part at $r=0$,
\begin{equation}\label{eq:trace-residue}
[R^{-K}]\,\mathcal C_{\Phi,5}(he)=\frac{1}{p}\Res_{r=0}\!\bigl(h(r)\,J_\Phi(r)^{-1}\,F'(r)\,F(r)^{K-1}\,\dd r\bigr).
\end{equation}
In particular $J_\Phi^{-1}\in\Z_p[\zeta_3][[r^3]]^\times$ is a $\mu_3$-invariant unit on the tame chart, so it does not affect the $r$-valuation, the $p$-valuation, or the $\mu_3$-character class of the residue integrand.
\end{lemma}

\begin{proposition}[Length-three Witt--Cartier pole estimate]\label{prop:pole-estimate}\label{prop:pole-estimate-star}
For every $\ell\in\{1,2,3\}$ and every $s$ with $\ell\le s\le 3$,
\begin{equation}\label{eq:star}
p^{s-\ell}\,\mathcal C_{\Phi,5}\!\bigl(r^{2+s(1-p)}\Z_p[\zeta_3][[r^3]]\,e\bigr)\subset r^{-1}\Z_p[\zeta_3][[r^3]]\,e\pmod{p^{4-\ell}}.
\end{equation}
\end{proposition}

\begin{proof}
By \Cref{lem:trace-residue-tame}, the principal-part coefficient in $R$ is
\begin{equation}\label{eq:residue-formula}
[R^{-K}]\,\mathcal C_{\Phi,5}(fe)=\frac{1}{p}\Res_{r=0}\!\bigl(f(r)\,J_\Phi(r)^{-1}\,F'(r)\,F(r)^{K-1}\,\dd r\bigr),
\end{equation}
where $J_\Phi^{-1}\in\Z_p[\zeta_3][[r^3]]^\times$ has $r$-valuation $0$, $p$-valuation $0$, and $\mu_3$-character class $0$, hence does not affect the estimates below; we suppress it in the rest of the proof and refer to \Cref{lem:trace-residue-tame} for justification. After identifying the target completed disk with $\operatorname{Spf}A$ and relabelling $R$ back to $r$, this is the $r^{-K}$ coefficient.

\medskip\noindent\textbf{Step 1: $\mathbb Z/3$-counting forces $K\equiv 1\pmod 3$.}\quad
$f\in r^{2+s(1-p)}\Z_p[\zeta_3][[r^3]]$ contributes $r$-exponent $\equiv 2+s(1-p)\equiv 2\pmod 3$ since $1-p\equiv 0\pmod 3$. The derivative
\[
F'(r)=p\,r^{p-1}+pD_1(r)+p^2D_2(r)+p^3D_3(r),\quad D_i(r)=\alpha_i(r^3)+3r^3\alpha_i'(r^3),
\]
has both terms in exponent class $\equiv 0\pmod 3$. The factor $F(r)^{K-1}$ has each factor of $F$ in class $\equiv 1\pmod 3$, so contributes $\equiv K-1\pmod 3$. A nonzero residue requires total exponent $-1\equiv 2\pmod 3$, giving $2+0+(K-1)\equiv 2\pmod 3$, hence $K\equiv 1\pmod 3$. For $K\ge 2$ this forces $K\ge 4$.

\medskip\noindent\textbf{Step 2: Valuation count.}\quad
Write $F(r)^{K-1}=r^{p(K-1)}(1+X)^{K-1}$ with $X=pr^{1-p}\alpha_1+p^2r^{1-p}\alpha_2+p^3r^{1-p}\alpha_3$. The binomial expansion selects $b\le K-1$ factors of $X$, each contributing minimal $r$-exponent $1-p$ and $p$-valuation at least $1$.

\emph{Derivative-high case} ($F'\to pr^{p-1}$): total $r$-exponent
\[
E_{\rm high}=2+s(1-p)+(p-1)+p(K-1)+b(1-p)+3h=1+s+b+p(K-b-s)+3h
\]
with $h\ge 0$. Residue condition $E_{\rm high}=-1$:
\[
(p-1)(s+b)=pK+2+3h,\qquad s+b-K=\frac{K+2+3h}{p-1}>0,
\]
hence $s+b\ge K+1$, so $b\ge K+1-s$. With $K\ge 4$, $b\ge 5-s$.

\emph{Derivative-low case} ($F'\to p^jD_j$, $j\ge 1$): similar computation gives
\[
(p-1)(s+b)=p(K-1)+3+3h,\qquad s+b-(K-1)=\frac{K+2+3h}{p-1}>0,
\]
hence $s+b\ge K$ and $b\ge K-s\ge 4-s$ when $K\ge 4$.

\emph{Control example} ($p=7$, $K=4$, $s=1$, $h=0$): low case requires $6(s+b)=24$, so $s+b=4$ and $b=3$, satisfying $b\le K-1=3$. High case requires $6(s+b)=30$, so $b=4>K-1=3$, hence does not occur in this minimal example.

\medskip\noindent\textbf{Step 3: Total valuation.}\quad
Before the external $1/p$, each integrand has $p$-valuation $\ge 1+b$ (the explicit $p$ from $F'$ plus at least $b$ from the $X$-expansion). After $1/p$, valuation $\ge b\ge 4-s$ (the minimum, from the low case). Multiplying by external $p^{s-\ell}$ in \eqref{eq:star} gives valuation $\ge (4-s)+(s-\ell)=4-\ell$. Thus every $r^{-K}$ coefficient with $K\ge 2$ vanishes mod $p^{4-\ell}$.

The suppressed $O(p^4)$ part of $F$ contributes valuation $\ge 4$ before $1/p$, hence $\ge 3$ after, already vanishing for all $\ell\in\{1,2,3\}$.

\medskip\noindent\textbf{Step 4: $K=1$.}\quad
$F^{K-1}=1$, and the residue lands in the $r^{-1}$ part of the principal part. Combining steps gives \eqref{eq:star}.
\end{proof}

\end{document}
